% !TEX root = ../main.tex
\lecture{1}{2026-09-17}{Day 4 Intro}

\section{Admin Stuff}
\begin{itemize}
    \item Uebungen
        \subitem PDF Leseaufträge OLAT \textbf{pruefrelevant} 
\end{itemize}

\section{Begriffsdefinition}
\begin{itemize}
    \item (\textit{bürgerliche}) ökonomie beschäftigt sich mit \textit{rationalen} Entscheidungen unter \textit{Knappheit} 
    \footnote{\textbf{Bemerkung:} in der Definition sind alle spezifischen Aspekten der kapitalistischen Gesellschaftsform (verallgemeinerter Markt, Privateigetum ProdMittel, Geld, Kapital usw.) abstrahiert und naturalisiert sie dadurch. Die gesellschaftlichen Verhältnissen/Umgebung, worin entscheidungen stattfinden werden nicht erläutert. Eine nicht legitime abstraktion.}
    \subitem \textit{Knappheit} weniger von einem Objekt vorhanden als wünschenswert
    \subitem \textit{Rational} konsistente, möglichst optimale Entscheidungen treffen anhand der Kosten und Nutzen (Konsequenzen) 
    \item Positive (deskriptive) Aussagen
    \item Negative (normative) Aussagen 
    \item "ökonomie ist immer politisch"
\end{itemize}

\section{Kosten und Nutzen}

Beispiel Skifahren. Kosten (Benefits) \textbf{ $C(x)$} und Nutzen (Benefits) \textbf{$B(x)$}. Wenn \textbf{$C(x) < B(x)$} führt Akteur Handlung aus, wenn \textbf{$C(x) > B(x)$} wird nicht ausgeführt.
\textbf{$C(x) = B(x)$} Indifferenz.

\textbf{Messeinheiten}
\begin{itemize}
    \item Kosten in Marktpreisen (Geldsumme)
    \item Nutzen in Zahlungsbereitschaft (Geldsumme)
        \footnote{hier auch wird von den konkreten aspekten der Aktivität und der ausführenden Person abstrahiert. Könnte innerhalb eines Warenarts legitim sein, aber qualitativ verschiedene Gebrauchswerte können nicht von ihren konkreten Nutzzwecken quantitativ abstrahiert werden}
\end{itemize}

\textbf{Generell rationale Alternativen}
\begin{itemize}
    \item Mengen der Auswahlalternativen $X = \{x_1, x_2, \dots, x_n\}$
    \item $B(x*) - C(x*) \ge B(x) - C(x)$ für alle $x \in X$
\end{itemize}

Entscheidend ist sind Nutzen einer alternative groesser als seine entsprechenden Kosten und Nettonutzen der anderen Option. Falls keine andere Option vorhanden, \textbf{nichts machen} voraussetzen.
\begin{theorem}
    \label{t_allgemein}
    $B(x_1) > C(x_1) + (B(x_2) - C(x_2))$ 
\end{theorem}
\pagebreak

\textbf{Definition Oportunitaetskosten} sind \textit{entgangene} Nettonutzen der zweitbesten Alternative und \textbf{müssen} bei rationalen Entscheidungen berücksichtigt werden.

\begin{theorem}
    \label{t_oppkosten}
    $B(x_2) - C(x_2)$
\end{theorem}


\subsection{Irrationalitäten}
\textbf{Definition versunkenen Kosten} (\textit{sunk cost fallacy}) sind Kosten die bei einer vergangenen Entscheidung gefallen sind, und die eine gegenwärtige oder künftige Entscheidung nicht beeinflüssen dürfen. 

\subsection{Egoismus und Altruismus}

\textbf{Definition} Homo Oeconomicus, trifft nur entscheidungen, die auf seinen eigenen Kosten fallen und ihren eigenen Nutzen maximieren. (\textit{Kann sich zugleich altruistisch und rational verhalten})



 